Residence near a four-year college produces a larger estimated return to schooling than ordinary least squares, consistent with attenuation in the OLS estimates. Among the 3,010 men in the NLS Young Men cohort, one additional year of schooling is associated with 5.3\% higher hourly wages in the baseline model (\textit{b} = 0.052, \textit{SE} = 0.003, \textit{p} < .001) and 7.8\% higher wages after controlling for race, region, urban residence, and experience (\textit{b} = 0.075, \textit{SE} = 0.003, \textit{p} < .001). Instrumenting schooling with an indicator for growing up near a four-year college (nearc4) raises the return to 14.1\% (\textit{b} = 0.132, \textit{SE} = 0.055, \textit{p} = .017, 95\% CI [0.024, 0.239]); the instrument is not weak (first-stage \textit{F}(1, 2994) = 13.26, \textit{p} < .001). If college proximity affects wages only through schooling, the gap between the OLS and 2SLS estimates implies that conventional estimates understate the return to schooling by about 6 percentage points; the wide 2SLS interval and the local-average interpretation of the IV estimand nonetheless caution against treating 14\% as a population-wide causal effect.
